\documentclass[11pt]{article}
\usepackage[T1]{fontenc}
\usepackage[utf8]{inputenc}
\usepackage{amsmath,amssymb,amsthm}
\usepackage[margin=1in]{geometry}
\usepackage{graphicx}
\usepackage{booktabs}
\usepackage{float}
\usepackage[colorlinks=true,linkcolor=blue,citecolor=blue]{hyperref}

\newtheorem{lemma}{Lemma}
\theoremstyle{remark}
\newtheorem{remark}{Remark}

\newcommand{\Salg}{S_{\mathrm{alg}}}

\title{Algebraic Entropy and Conditional Mutual Information in a Tiny
Gauge-Invariant Truncated Hilbert Space:\\
\large A Reproducible Toy-Model Study with Effective Mixing
Hamiltonians}
\author{Lluis Eriksson\\
\small Independent Researcher\\
\small \texttt{lluiseriksson@gmail.com}}
\date{January 2026 (v1); July 2026 (v2)}

\begin{document}
\maketitle

\begin{abstract}
We present a reproducible pipeline to compute region algebraic entropies
and conditional mutual informations (CMI) in a tiny truncated Hilbert
space (here $\dim = 8$) indexed by discrete fusion-like descriptors
$\mathrm{desc} = (x, \mu)$ on $L = 4$ cells. To generate nontrivial
ground states within the descriptor-labeled subspace, we introduce an
effective Hermitian mixing Hamiltonian based on a weighted
$k$-nearest-neighbor (kNN) graph Laplacian over configuration labels.
Across a parameter sweep, we identify a strong-mixing regime where the
participation ratio approaches $\dim$ (consistent with
Laplacian-dominated ground states on connected graphs) and algebraic CMI
diagnostics become extremely small (down to $10^{-6}$ and below) for the
chosen algebraic factorization, while region algebraic entropies remain
$O(1)$ and exhibit near-quantized values $\approx n\log 2$. We stress
that the mixing term is an ansatz used to probe information-theoretic
diagnostics and is not claimed to coincide with a Kogut--Susskind
plaquette operator. \emph{v2 (no v1 number is changed):} the
reproducibility gap of v1 is repaired --- v1's pipeline loaded an
unshipped basis file (\texttt{descs.pkl}) that was never specified, so
the v1 dataset was not regenerable from the paper; v2 prints a canonical
self-contained basis (Appendix~\ref{app:canonical}) whose pipeline
reproduces every structural finding, keeping v1's Table~\ref{tab:models}
as an archival dataset; two empirical observations are upgraded to
proved statements --- the descriptor-to-key map is injective, making
$\Salg$ a genuine von Neumann entropy of a sector (center-type)
decomposition (Lemma~\ref{lem:inj}), and in the strong-mixing limit the
ground state converges to the uniform superposition where the
quantization $\Salg = n\log 2$ and the vanishing of both CMIs are
\emph{exact} (Lemma~\ref{lem:uniform}); the additional experiments
recommended in v1 (Haar baseline, $k_{\mathrm{nn}}$ ablations, finer
$t_{\mathrm{mix}}$ grid) are executed --- the Haar median of
$I_{\mathrm{sum}}$ is $0.39$, five to six orders of magnitude above the
strong-mixing point, so the small-CMI regime is nontrivial; and the
verification suite is fully self-contained (no Drive dependencies).
\end{abstract}

\section{Scope, claims, and contributions}

This note is intentionally narrow in scope. It does \emph{not} attempt a
physical construction of lattice Yang--Mills dynamics, nor any
continuum/thermodynamic scaling analysis. Instead, it addresses a
simpler, operational question:

\emph{Given a fixed small gauge-invariant truncated Hilbert space and a
specified algebraic factorization derived from discrete labels $(x,
\mu)$, can we generate nontrivial ground states and quantitatively
diagnose algebraic entanglement and conditional independence via
algebraic entropy and algebraic CMI?}

Concretely, we contribute:
\begin{enumerate}
\item A one-shot reproducible pipeline to compute algebraic entropies
$\Salg(R)$ and algebraic CMI diagnostics $I(A : C \mid B)$ from a
gauge-invariant basis indexed by $(x, \mu)$ (self-contained as of v2).
\item A simple, tunable effective Hamiltonian family $H = H_E +
t_{\mathrm{mix}} L$ acting entirely within the gauge-invariant subspace
(Section~5).
\item A numerical sweep showing that in a strong-mixing regime the
ground state becomes close to a uniform superposition ($\mathrm{PR}
\approx \dim$), while the resulting algebraic CMI diagnostics are
near-zero within numerical precision for the chosen factorization
(exactly zero in the uniform limit, Lemma~\ref{lem:uniform}).
\end{enumerate}

\textbf{Careful interpretation.} When we refer to ``approximately
Markovian'' behavior, we mean an \emph{operational} statement: the
computed algebraic CMI values are very small with respect to the
specific algebraic partition induced by $(x, \mu)$ and the implemented
definition of $\Salg$. We do not claim a general quantum Markov property
in the usual tensor-product sense, nor any universality beyond this toy
setting.

\section{Truncated gauge-invariant Hilbert space and labels}

We work in a fixed truncated Hilbert space with orthonormal basis
$\{|i\rangle\}_{i=1}^{\dim}$, where each basis vector is indexed by a
descriptor $\mathrm{desc} = (x, \mu)$:
\begin{itemize}
\item $x = (x_1, x_2, x_3)$ encodes boundary/cut labels between adjacent
cells. Concretely, each $x_s = (x_s^{\mathrm{top}}, x_s^{\mathrm{bot}})
\in \{0, 1\}^2$ stores two binary cut labels.
\item $\mu = (\mu_1, \mu_2, \mu_3, \mu_4)$ encodes local discrete
degrees of freedom per cell; each $\mu_p$ is a (possibly multi-bit)
tuple.
\end{itemize}
We take $L = 4$ cells and, for the dataset studied here, $\dim = 8$.

\textbf{Terminology.} We refer to the basis as ``gauge-invariant'' in
the operational sense that it is intended to represent
constraint-compatible (Gauss-law-like) configurations in a reduced
description; however, in this note we do not derive $(x, \mu)$ from an
explicit $SU(2)$ Kogut--Susskind truncation.

\textbf{Basis provenance (v2).} The v1 dataset was generated from a
basis file \texttt{descs.pkl} loaded from persistent storage; that file
was never printed or specified in the paper, so the v1 numbers in
Table~\ref{tab:models} are archival (kept unchanged) but not regenerable
from the paper alone. v2 fixes this by specifying a \emph{canonical}
self-contained basis (Appendix~\ref{app:canonical}) with the same label
structure and $\dim = 8$; all structural claims of this paper are
regenerated from it by the verification suite.

\textbf{Logarithms.} All entropies use natural logarithms, so the unit
is nats. Thus $\log 2 \approx 0.6931$.

\section{Regions and target information-theoretic quantities}

We consider contiguous regions $R \subset \{1, 2, 3, 4\}$ and compute
algebraic entropies $\Salg(R)$. We define two conditional mutual
informations:
\[
I_{w=1} = I(A : C \mid B), \quad A = \{1\},\ B = \{2\},\ C = \{3\},
\]
\[
I_{w=2} = I(A : C \mid B), \quad A = \{1\},\ B = \{2, 3\},\ C = \{4\}.
\]
We also report $I_{\mathrm{sum}} = I_{w=1} + I_{w=2}$.

\textbf{Purity sanity check.} For a pure state $|\psi\rangle$ on the
full system, we enforce $\Salg(1..4) \approx 0$ as a numerical check.

\section{Definition of algebraic entropy and algebraic CMI (as
implemented)}

Let $|\psi\rangle = \sum_{i=1}^{\dim} \psi_i |i\rangle$ be a normalized
pure state in the gauge-invariant basis. For a region $R = [i..j]$, we
define a discrete boundary label $\alpha = \alpha(\mathrm{desc}; i, j)$
(a tuple extracted from $x$ at cut boundaries) which indexes
superselection sectors for the algebraic factorization.

Within each sector $\alpha$, basis states are grouped into a region key
$k_R$ and complement key $k_C$ (deterministic functions of $(x, \mu)$,
defined identically to the accompanying code). We build a complex matrix
$M_\alpha$ whose entries sum amplitudes of basis states mapping to the
same $(k_R, k_C)$ pair:
\[
(M_\alpha)_{r,c} = \sum_{\mathrm{desc} \mapsto (\alpha, r, c)}
\psi(\mathrm{desc}).
\]
Let $\{s_{\alpha,\ell}\}_\ell$ be the singular values of $M_\alpha$, and
define weights $w_{\alpha,\ell} = s^2_{\alpha,\ell}$. The algebraic
entropy is
\[
\Salg(R) = -\sum_\alpha \sum_{\ell:\, w_{\alpha,\ell} > \varepsilon}
w_{\alpha,\ell} \log w_{\alpha,\ell},
\]
where a small cutoff $\varepsilon$ discards numerical zero modes only.

\textbf{Normalization check.} For normalized $|\psi\rangle$, we
numerically verify $\sum_{\alpha,\ell} w_{\alpha,\ell} = 1$ within a
fixed tolerance.

\begin{lemma}[Injectivity: $\Salg$ is a von Neumann entropy (v2)]
\label{lem:inj}
Suppose the map $\mathrm{desc} \mapsto (\alpha, k_R, k_C)$ is injective
on the basis (as holds for the canonical basis of
Appendix~\ref{app:canonical}, verified for all six regions used, and as
is certified in any run by the normalization check $\sum w = 1$: if two
basis states shared a key triple, cross terms would generically break
$\sum_{\alpha,\ell} w_{\alpha,\ell} = 1$). Then $(M_\alpha)_{r,c}$ is a
mere re-indexing of the amplitude vector, the map $\psi \mapsto
\bigoplus_\alpha M_\alpha$ is an isometry onto $\bigoplus_\alpha
\mathcal{H}_R^\alpha \otimes \mathcal{H}_C^\alpha$, and
\[
\Salg(R) = S\bigl(\rho_R^{\mathrm{alg}}\bigr), \qquad
\rho_R^{\mathrm{alg}} = \bigoplus_\alpha \mathrm{Tr}_{C}\,
M_\alpha M_\alpha^\dagger,
\]
i.e.\ $\Salg$ is the genuine von Neumann entropy of the reduced state on
a direct-sum (superselection-sector) decomposition --- the standard
center-type structure for gauge-theory entanglement \cite{CHR}. It
therefore includes the usual sector-classical (Shannon) plus
sector-quantum contributions.
\end{lemma}

\begin{proof}
Injectivity means every basis amplitude $\psi(\mathrm{desc})$ lands in
exactly one entry of one $M_\alpha$, so $\sum_{\alpha}
\|M_\alpha\|_F^2 = \|\psi\|^2 = 1$ and the map is an isometry onto its
image. The weights $w_{\alpha,\ell}$ are then the eigenvalues of
$\bigoplus_\alpha M_\alpha M_\alpha^\dagger$, which is exactly the
reduced density matrix of the image state on the $R$-legs of the
direct sum.
\end{proof}

\textbf{Algebraic CMI diagnostics.} We compute
\[
I(A : C \mid B) = \Salg(AB) + \Salg(BC) - \Salg(B) - \Salg(ABC).
\]
In the usual tensor-product setting, CMI is nonnegative by strong
subadditivity. In our setting $\Salg$ is defined by an explicit
sector/block coarse-graining, so nonnegativity is an empirical
diagnostic rather than a guaranteed theorem
(Lemma~\ref{lem:inj} makes each individual $\Salg(R)$ a genuine vN
entropy, but the four regions' sector decompositions are not derived
from one consistent assignment of commuting subalgebras, so SSA does not
automatically apply). In our sweeps we observe $I(A : C \mid B) \ge
-10^{-12}$; values below this threshold would indicate either numerical
issues or a genuine violation of strong subadditivity for the chosen
coarse-graining. For log-scale plots one may use the clipped version
$I_{\mathrm{clip}} = \max(I, 10^{-16})$.

\section{Effective Hamiltonian family}

We study an effective Hamiltonian family acting within the
gauge-invariant subspace: $H = H_E + t_{\mathrm{mix}} L$.

\subsection{Diagonal ``electric proxy'' term $H_E$}

We define a diagonal proxy $H_E$ as a sum of local contributions derived
from the binary label bits. Each bit $b \in \{0, 1\}$ is mapped to a
cost inspired by $j(j+1)$ for $j \in \{0, \frac12\}$: $\mathrm{cost}(b)
= 0$ for $b = 0$ and $\frac34$ for $b = 1$. Then $H_E =
\mathrm{diag}(E(\mathrm{desc}))$, where $E(\mathrm{desc})$ sums costs
over all bits in $(x, \mu)$, with an overall prefactor $\frac{g^2}{2}$.

\subsection{Mixing term via a weighted kNN graph Laplacian}

Let $b(\mathrm{desc}) \in \{0, 1\}^m$ be the flattened bitstring for
$(x, \mu)$. Define pairwise distances by Hamming distance $d(i, j)$
between bitstrings. We build a weighted kNN graph on the $\dim$
configurations: for each node $i$ we connect to its $k$ nearest
neighbors, with symmetric edge weights $A_{ij} = \exp(-\alpha_w d(i,
j))$, and define the graph Laplacian $L = D - A$, $D_{ii} = \sum_j
A_{ij}$. This yields a real symmetric (hence Hermitian) mixing term. The
design intent is to generate controllable superpositions in the
gauge-invariant basis. No claim is made that $L$ matches a
Kogut--Susskind magnetic plaquette term \cite{KS}.

\begin{lemma}[Strong-mixing limit (v2)]
\label{lem:uniform}
Fix $\alpha_w = 0$ and $k$ such that the kNN graph is connected. As
$t_{\mathrm{mix}} \to \infty$, the ground state of $H = H_E +
t_{\mathrm{mix}} L$ converges to the uniform superposition
$|u\rangle = \dim^{-1/2}\sum_i |i\rangle$ (the Perron--Frobenius ground
state of $L$), so $\mathrm{PR} \to \dim$. At $|u\rangle$, for the
canonical basis of Appendix~\ref{app:canonical}, the block combinatorics
give \emph{exactly}
\[
\Salg(1..2) = \Salg(1..3) = \Salg(2..4) = \log 2, \qquad
\Salg(2) = \Salg(2..3) = 2\log 2, \qquad \Salg(1..4) = 0,
\]
and hence $I_{w=1} = I_{w=2} = 0$ exactly: the near-quantization and
near-zero CMI observed in the sweep are the finite-$t_{\mathrm{mix}}$
shadow of an exact property of the uniform limit.
\end{lemma}

\begin{proof}
For a connected graph, $L \succeq 0$ with a unique zero mode $\propto
\mathbf{1}$; degenerate perturbation theory in $H_E/t_{\mathrm{mix}}$
gives $|\psi_0\rangle = |u\rangle + O(1/t_{\mathrm{mix}})$. At
$|u\rangle$ every $M_\alpha$ has all entries $\dim^{-1/2}$; a $d_R
\times d_C$ all-equal matrix has the single nonzero singular value
$\sqrt{d_R d_C}/\sqrt{\dim}$, so each sector contributes one weight
$w_\alpha = d_R d_C / \dim$ and $\Salg = -\sum_\alpha w_\alpha \log
w_\alpha$ is the Shannon entropy of the sector distribution. For the
canonical basis the sector counts are $(2, 4, 4, 2, 2, 1)$ sectors of
equal weight for the six regions, giving the displayed values; the two
CMI combinations cancel exactly. The suite verifies all six values and
both CMIs to $2\times10^{-16}$.
\end{proof}

\section{Numerical procedure and reported metrics}

For each parameter pair $(\alpha_w, t_{\mathrm{mix}})$ on a fixed grid:
(1) construct $H$, diagonalize it, and take the ground state
$|\psi_0\rangle$; (2) compute region entropies $\Salg(R)$, CMI
diagnostics $I_{w=1}$, $I_{w=2}$, and $I_{\mathrm{sum}}$; (3) compute
the participation ratio $\mathrm{PR}(\psi_0) = 1/\sum_i |\psi_{0,i}|^4$
and the spectral gap proxy $\Delta = E_1 - E_0$, $\xi_{\mathrm{spec}} =
\Delta^{-1}$ ($\Delta > 0$).

\textbf{Interpretation of $\xi_{\mathrm{spec}}$.} We emphasize that
$\xi_{\mathrm{spec}} = \Delta^{-1}$ is a finite-dimensional spectral
proxy whose value depends on the overall energy scale, in particular on
$t_{\mathrm{mix}}$. No claim is made that it corresponds to a physical
correlation length in any continuum or infinite-volume limit.

\section{Artifacts}

The v1 one-shot run generated \texttt{sweep.csv}/\texttt{sweep.json},
the three figures, selected model dumps and an Overleaf bundle (see the
v1 PDF for the full Colab script). \emph{v2:} the verification suite
(\texttt{verification/2601-0115/verify\_2601\_0115.py} in the series
repository) is fully self-contained --- no Drive mount, no
\texttt{descs.pkl} --- and regenerates the sweep, the lemma checks, and
the additional experiments below from the canonical basis.

\section{Results}

Figures~\ref{fig:pr}--\ref{fig:s23} show the v1 sweep (archival).
Table~\ref{tab:models} reports the two automatically selected models of
the v1 run: maximum PR, and minimum $I_{\mathrm{sum}}$ under $\mathrm{PR}
\ge 2$; the strong-mixing point dominates and yields PR close to $\dim$
with very small CMI.

\begin{figure}[h!]
\centering
\includegraphics[width=0.62\linewidth]{fig1.png}
\caption{Participation ratio (PR) versus total algebraic CMI
$I_{\mathrm{sum}} = I_{w=1} + I_{w=2}$ across the v1 parameter sweep
(log scale on CMI).}
\label{fig:pr}
\end{figure}

\begin{figure}[h!]
\centering
\includegraphics[width=0.62\linewidth]{fig2.png}
\caption{Spectral gap proxy $\Delta$ versus total algebraic CMI across
the v1 sweep (log scale on CMI).}
\label{fig:gap}
\end{figure}

\begin{figure}[h!]
\centering
\includegraphics[width=0.5\linewidth]{fig3.png}
\caption{Entropy-structure scatter: $\Salg(2..3)$ versus
$\Salg(1..3)$ (v1 sweep).}
\label{fig:s23}
\end{figure}

\begin{table}[H]
\centering
\caption{Automatically selected models from the v1 sweep (archival; as
reported in \texttt{paper\_summary.txt}).}
\label{tab:models}
\begin{tabular}{l c c}
\toprule
Quantity & Best-by-PR & Best-by-CMI ($\mathrm{PR} \ge 2$) \\
\midrule
$\alpha_w$ & 0.0 & 0.0 \\
$t_{\mathrm{mix}}$ & 100.0 & 100.0 \\
$k_{\mathrm{nn}}$ & 3 & 3 \\
edges & 12 & 12 \\
$E_0$ & 2.6179476 & 2.6179476 \\
$\Delta$ & 199.4735852 & 199.4735852 \\
$\xi_{\mathrm{spec}} = \Delta^{-1}$ & 0.0050132 & 0.0050132 \\
PR & 7.9989229 & 7.9989229 \\
$\Salg(1..2)$ & 0.6931399 & 0.6931399 \\
$\Salg(2..3)$ & 1.3862379 & 1.3862379 \\
$\Salg(2)$ & 1.3862571 & 1.3862571 \\
$\Salg(1..3)$ & 0.6931190 & 0.6931190 \\
$\Salg(2..4)$ & 0.6931190 & 0.6931190 \\
$\Salg(1..4)$ & $2.2\times10^{-16}$ & $2.2\times10^{-16}$ \\
$I_{w=1}$ & $1.76\times10^{-6}$ & $1.76\times10^{-6}$ \\
$I_{w=2}$ & $1.39\times10^{-11}$ & $1.39\times10^{-11}$ \\
$I_{\mathrm{sum}}$ & $1.76\times10^{-6}$ & $1.76\times10^{-6}$ \\
\bottomrule
\end{tabular}
\end{table}

\textbf{Canonical-basis regeneration (v2).} With the canonical basis,
the strong-mixing point $(\alpha_w, t_{\mathrm{mix}}) = (0, 100)$ gives
$\mathrm{PR} = 7.9994$ and $I_{\mathrm{sum}} = 4.4\times10^{-7}$, with
every $\Salg$ within $10^{-3}$ of its exact uniform-limit value ---
structurally identical to the archival Table~\ref{tab:models} (whose
entropy pattern $(\log 2, 2\log 2, 2\log 2, \log 2, \log 2)$ matches the
canonical uniform limit of Lemma~\ref{lem:uniform} to four decimals),
confirming that the findings do not depend on the lost v1 basis file.

\textbf{Executed additional experiments (v2; recommended in v1
\S9.4).}
\emph{(i) Haar baseline:} over $2000$ Haar-random pure states in
$\mathbb{C}^8$, the empirical $I_{\mathrm{sum}}$ has median $0.389$
(5th percentile $0.158$) --- five to six orders of magnitude above the
strong-mixing value, so the small-CMI regime is a genuinely nontrivial
region of state space, not a generic feature of the diagnostic.
Empirical SSA holds for all Haar samples. \emph{(ii) $k_{\mathrm{nn}}$
ablation:} for $k \in \{1, 2, 5\}$ the kNN graph remains connected for
the canonical basis and the strong-mixing conclusions persist
($\mathrm{PR} > 7.9$, $I_{\mathrm{sum}} \le 10^{-4}$), with CMI
decreasing as connectivity grows. \emph{(iii) Finer $t_{\mathrm{mix}}$
grid:} $\Salg(2..3)$ rises monotonically from $0$ (diagonal regime,
$t = 0$) to $\approx 2\log 2$ across $t \in \{0, 0.1, 0.3, 1, 3, 10\}$,
resolving the transition to the Laplacian-dominated regime.

\section{Discussion}

\subsection{Why $\mathrm{PR} \approx \dim$ and $\mathrm{CMI} \approx 0$
occur in the strong-mixing regime}

At $\alpha_w = 0$ the kNN adjacency weights satisfy $A_{ij} = 1$ on
selected edges, so $L$ becomes the combinatorial graph Laplacian of a
regular graph (in the reported run, degree 3). For a connected graph,
$L$ has a lowest eigenvalue $0$ with eigenvector proportional to the
all-ones vector $\mathbf{1}$. Therefore, when $t_{\mathrm{mix}}$
dominates the diagonal proxy term $H_E$, the ground state approaches a
near-uniform superposition, yielding PR close to $\dim$. In this regime
the near-quantized values $\Salg(R) \approx n\log2$ and the very small
algebraic CMI values are the finite-$t$ shadow of the exact uniform-limit
statement of Lemma~\ref{lem:uniform}. Operationally, this indicates
near-conditional-independence \emph{with respect to that algebraic
factorization}.

\subsection{What this does (and does not) say about ``quantum Markov''
structure}

In standard quantum information theory, exact saturation of strong
subadditivity ($I(A : C \mid B) = 0$) implies a quantum Markov structure
and is characterized by a recovery map \cite{HJPW}. Our setting differs
because: entropies are computed via a sector-decomposed algebraic
prescription tied to $(x, \mu)$ labels; the system is extremely small
($\dim = 8$); and the Hamiltonian mixing term is an effective graph
Laplacian ansatz. Thus, small values of the computed algebraic CMI
should be read as a diagnostic of conditional-independence in this toy,
label-induced algebraic setting, rather than as evidence for a
physically universal Markov property.

\subsection{Limitations}

\textbf{No scaling.} We do not study how observables behave as $\dim$ or
$L$ increase. \textbf{Hamiltonian not physical.} The mixing term is not
derived from a Kogut--Susskind magnetic term; it is a controllable
ansatz. \textbf{Spectral proxy.} $\xi_{\mathrm{spec}} = \Delta^{-1}$ is
scale-dependent and should not be overinterpreted.

\subsection{Series positioning (v2)}

Within the series, this note is the algebraic-toy end of the gauge/CMI
block: 2601.0111 (v2) computes tensor-product CMI and Petz recovery in
an exact $\mathbb{Z}_2$ Gauss sector; 2601.0043 (v2) provides the
embedding demonstrations for algebraic-vs-tensor comparisons; the
center-type sector structure made explicit by Lemma~\ref{lem:inj} is the
finite toy version of the split/sector discipline of the Type III
interface 2601.0065 (v2). The Haar baseline executed here is the
toy-scale analogue of the null-model controls used across the series'
numerical notes.

\section{Numerical precision and stability}

Occasional values of CMI at the level of $10^{-16}$ (including slight
negatives) are interpreted as floating-point round-off in double
precision accumulated over many linear-algebra operations. For log-scale
plots, we use $\max(I_{\mathrm{sum}}, 10^{-16})$, and when reporting
aggregated metrics we may optionally use $I_{\mathrm{clip}} = \max(I,
0)$.

\section{Reproducibility}

The v1 Colab cell (typeset in full in the v1 PDF, Appendix B) produced
timestamped run folders under Drive and required a pre-existing
\texttt{descs.pkl}. \emph{v2:} the self-contained verification suite
replaces this dependency entirely; a single \texttt{python3
verify\_2601\_0115.py} regenerates every structural claim, the lemma
checks, the sweep, and the executed additional experiments from the
canonical basis of Appendix~\ref{app:canonical}.

\section*{Note added in v2}
No v1 number is changed; Table~\ref{tab:models} and the three figures
are kept as the archival v1 dataset. v2: (1) repairs the reproducibility
gap --- v1's pipeline loaded an unshipped, never-specified
\texttt{descs.pkl}, so the v1 dataset was not regenerable from the paper
alone; v2 prints a canonical self-contained basis
(Appendix~\ref{app:canonical}) and demotes the v1 dataset to archival
status while regenerating every structural finding from the canonical
basis (strong-mixing point: $\mathrm{PR} = 7.9994$, $I_{\mathrm{sum}} =
4.4\times10^{-7}$, entropy pattern identical); (2) proves the
injectivity lemma (Lemma~\ref{lem:inj}): $\Salg$ is a genuine von
Neumann entropy of a center-type sector decomposition \cite{CHR}, and
clarifies why SSA nonetheless remains empirical for the CMI
combinations; (3) proves the strong-mixing lemma
(Lemma~\ref{lem:uniform}): the quantization $\Salg = n\log2$ and
$I_{w=1} = I_{w=2} = 0$ are exact in the uniform limit (verified to
$2\times10^{-16}$); (4) executes the three additional experiments
recommended in v1 \S9.4 (Haar baseline, $k_{\mathrm{nn}}$ ablation,
finer $t_{\mathrm{mix}}$ grid); (5) adds series positioning (v1 cited no
companion papers) and the self-contained verification suite.

\appendix

\section{Canonical self-contained basis (v2)}
\label{app:canonical}

The canonical basis has $\dim = 8$ states labeled by three cut bits
$(b_1, b_2, b_3) \in \{0,1\}^3$:
\[
x_s = (b_s, b_s) \ \ (s = 1, 2, 3), \qquad
\mu_1 = (b_1), \quad \mu_2 = (b_1 \oplus b_2), \quad
\mu_3 = (b_2 \oplus b_3), \quad \mu_4 = (b_3).
\]
The two cut bits per boundary are equal (a constraint-compatibility
condition) and the cell labels $\mu_p$ are parity-matched to the
adjacent cuts (a Gauss-law-like consistency), so every descriptor is
determined by the cut configuration. The flattened bitstring has $m =
10$ bits (6 from $x$, 4 from $\mu$). The map $\mathrm{desc} \mapsto
(\alpha, k_R, k_C)$ is injective for all six regions used
(Lemma~\ref{lem:inj}; verified exhaustively), the kNN graph at $k = 3$,
$\alpha_w = 0$ is connected with 12 edges (3-regular, as in the v1 run),
and the uniform-limit sector counts are $(2, 4, 4, 2, 2, 1)$ for the
regions $(1..2),\ (2..3),\ (2),\ (1..3),\ (2..4),\ (1..4)$, giving the
exact entropies of Lemma~\ref{lem:uniform}.

\section{Minimal pseudocode overview}

The workflow is: (1) build the canonical descriptors (or, for the
archival v1 dataset, load \texttt{descs.pkl}); (2) build bitstrings
$b(\mathrm{desc}_i)$ and the Hamming distance matrix $D_{ij}$; (3) for
each $(\alpha_w, t_{\mathrm{mix}})$: build the kNN adjacency $A$ and
Laplacian $L$; build $H = H_E + t_{\mathrm{mix}} L$; diagonalize and
take the normalized ground state; compute $\Salg(R)$ via sectored SVD;
compute $I_{w=1}$, $I_{w=2}$, PR, $\Delta$; save to CSV/JSON. The full
v1 Colab script remains typeset in the v1 PDF; the v2 suite supersedes
it with a dependency-free implementation.

\section{Notes on related literature (brief)}

The present work is primarily a reproducible toy-model pipeline. For
context, entanglement in gauge theories requires care due to constraints
and centers (see e.g.\ \cite{CHR}). The Kogut--Susskind Hamiltonian
formulation provides the standard lattice gauge Hamiltonian setting
\cite{KS}. In standard quantum information, the condition $I(A : C \mid
B) = 0$ characterizes quantum Markov states and admits structural
theorems \cite{HJPW}.

\begin{thebibliography}{9}
\bibitem{KS} J. Kogut and L. Susskind, ``Hamiltonian formulation of
Wilson's lattice gauge theories,'' Phys. Rev. D \textbf{11} (1975)
395--408.
\bibitem{CHR} H. Casini, M. Huerta, and J. A. Rosabal, ``Remarks on
entanglement entropy for gauge fields,'' Phys. Rev. D \textbf{89} (2014)
085012.
\bibitem{HJPW} P. Hayden, R. Jozsa, D. Petz, and A. Winter, ``Structure
of states which satisfy strong subadditivity of quantum entropy with
equality,'' Commun. Math. Phys. \textbf{246} (2004) 359--374.
\bibitem{S0111} L. Eriksson, ai.viXra:2601.0111 (2026; v2 2026).
\bibitem{S0043} L. Eriksson, ai.viXra:2601.0043 (2026; v2 2026).
\bibitem{S0065} L. Eriksson, ai.viXra:2601.0065 (2026; v2 2026).
\bibitem{S0050} L. Eriksson, ai.viXra:2601.0050 (2026; v2 2026).
\end{thebibliography}

\end{document}
